\documentclass[11pt,letterpaper]{article}

\usepackage[spanish,es-noshorthands]{babel}
\usepackage{fontspec}
\setmainfont{Source Serif Pro}
\usepackage{microtype}
\usepackage{geometry}
\usepackage{enumitem}
\usepackage{xcolor}
\usepackage{csquotes}
\usepackage{float}
\usepackage{tikz}
\usetikzlibrary{arrows.meta,positioning,fit,backgrounds}
\usepackage{pgfplots}
\pgfplotsset{compat=1.18}
\usepackage{xurl}
\usepackage[hidelinks]{hyperref}
\usepackage[backend=biber,style=apa,sorting=nyt,maxcitenames=2]{biblatex}
\DeclareLanguageMapping{spanish}{spanish-apa}

\geometry{margin=2.3cm}
\hbadness=10000
\setlength{\emergencystretch}{1em}
\setlength{\parindent}{0pt}
\setlength{\parskip}{0.5em}
\setlist{nosep}
\addbibresource{referencias.bib}
\AtBeginBibliography{\sloppy\small}
\setlength{\bibitemsep}{0.3em}

% Colores: verde para Apple Silicon, terracota para el rendimiento multinúcleo,
% azul oscuro para ejes y rótulos y gris para la competencia.
\definecolor{cApple}{RGB}{36,122,82}
\definecolor{cMulti}{RGB}{186,86,30}
\definecolor{cEje}{RGB}{26,58,102}
\definecolor{cGris}{RGB}{120,116,110}

\pgfplotsset{
  foro/.style={
    font=\small,
    axis line style={cEje},
    tick style={cEje},
    every tick label/.append style={font=\small, text=cEje},
    label style={font=\small, text=cEje},
    /pgf/number format/use comma,
    /pgf/number format/1000 sep={\,},
  },
}

\begin{document}

% --- Portada -----------------------------------------------------------------
\begin{titlepage}
\centering
\vspace*{3.2cm}
{\large\scshape Gestión del Desarrollo Tecnológico:\\
Estrategias y Análisis de Portfolio\par}
\vspace{0.5cm}
{\normalsize Módulo 6. Foro 6.2: Ejemplos de innovaciones tecnológicas\\
que han contribuido al aumento de la productividad\par}
\vspace{2.6cm}
{\color{cEje}\rule{0.55\textwidth}{0.8pt}\par}
\vspace{0.7cm}
{\LARGE\bfseries De M1 a M6\par}
\vspace{0.5cm}
{\large Cómo Apple Silicon cambió lo que una computadora\\
puede hacer con cada vatio\par}
\vspace{0.7cm}
{\color{cEje}\rule{0.55\textwidth}{0.8pt}\par}
\vfill
{\large Carlos Luis Rojas Aragonés\par}
\vspace{0.3cm}
{\normalsize Chief Technology Officer Program\par}
\vspace{0.3cm}
{\normalsize 10 de octubre de 2026\par}
\vspace*{1.5cm}
\end{titlepage}

% --- Cuerpo ------------------------------------------------------------------
Este comentario resume un artículo que publiqué en mi blog
\parencite{rojas2026}, a partir de lo que estudié en este programa.

\section*{1. La innovación: el chip, la memoria y el sistema como un solo diseño}

En junio de 2020 Apple anunció que el Mac dejaría los procesadores de Intel
por chips propios \parencite{apple2020transicion}. El M1, de noviembre de ese
año, integró CPU, GPU y aceleradores en un sistema en un chip, con memoria
unificada en el mismo paquete \parencite{apple2020m1}. La innovación no fue
solo el procesador: Apple diseñó chip, sistema operativo y producto a la vez,
y preparó la transición del software con Rosetta~2 y los binarios Universal~2.
Aun así, el ecosistema tardó en ponerse al día: Docker Desktop no tuvo
soporte estable hasta abril de 2021 \parencite{docker2021}.

\section*{2. Efectos en el rendimiento}

El MacBook Air con M1 eliminó el ventilador y anunció hasta 18 horas de
reproducción de video: el silencio dejó de depender de una carga ligera. Seis
generaciones después, el rendimiento de un núcleo creció 1,9 veces y el
multinúcleo 3,1 veces, en parte porque se pasó de 8 a 12 núcleos
\parencite{timm2026}. Lo muestra la Figura~\ref{fig:evolucion}.

\begin{figure}[H]
\centering
\begin{tikzpicture}
\begin{axis}[foro,
    width=12cm, height=6.2cm,
    symbolic x coords={M1,M2,M3,M4,M5,M6}, xtick=data,
    ymin=0.8, ymax=3.4, ytick={1,1.5,2,2.5,3},
    yticklabel={\pgfmathprintnumber{\tick}×},
    ylabel={Rendimiento relativo a M1},
    ymajorgrids, grid style={cGris!25},
    axis x line*=bottom, axis y line*=left,
    enlarge x limits=0.06, clip=false]
  \addplot[cMulti, very thick, mark=*, mark size=2.2pt]
    table[x=chip, y expr=\thisrow{multi}/1816] {
      chip multi
      M1 1816
      M2 2183
      M3 2558
      M4 3600
      M5 3957
      M6 5674
    } node[pos=1, anchor=west, xshift=4pt, font=\small\bfseries, text=cMulti]
      {3,1×};
  \addplot[cApple, very thick, mark=*, mark size=2.2pt]
    table[x=chip, y expr=\thisrow{mono}/428] {
      chip mono
      M1 428
      M2 461
      M3 545
      M4 644
      M5 733
      M6 823
    } node[pos=1, anchor=west, xshift=4pt, font=\small\bfseries, text=cApple]
      {1,9×};
  \node[font=\small, text=cMulti, anchor=west] at (axis cs:M2,2.75)
    {Multinúcleo: de 1\,816 a 5\,674 puntos};
  \node[font=\small, text=cApple, anchor=west] at (axis cs:M2,2.45)
    {Un núcleo: de 428 a 823 puntos};
\end{axis}
\end{tikzpicture}
\caption{Evolución de los chips base de Apple Silicon en Cinebench 2026.1,
normalizada respecto a M1. Son agregados de la comunidad con configuraciones y
refrigeración distintas, y M6 tiene una sola observación \parencite{timm2026}.}
\label{fig:evolucion}
\end{figure}

Lo más relevante, sin embargo, es el trabajo por vatio. En las pruebas de
\textcite{notebookcheck2026}, que miden la potencia de todo el sistema, el M6
supera en eficiencia multinúcleo a los equipos de Intel, AMD y Qualcomm
evaluados (Figura~\ref{fig:eficiencia}). No es invencible: el MacBook Air
con M4 rinde más puntos por vatio y el Snapdragon X2 Elite Extreme logra más
rendimiento multinúcleo absoluto.

\begin{figure}[H]
\centering
\begin{tikzpicture}
\begin{axis}[foro,
    width=12cm, height=6.2cm,
    xbar, bar width=9pt,
    xmin=0, xmax=45,
    symbolic y coords={AMD,Intel,Snapdragon,M5,M6,MacBook Air M4},
    ytick={AMD,Intel,Snapdragon,M5,M6,MacBook Air M4},
    bar shift=0pt,
    xlabel={Puntos multinúcleo por vatio (Cinebench 2024)},
    xmajorgrids, grid style={cGris!25},
    axis x line*=bottom, axis y line*=left,
    nodes near coords, every node near coord/.append style={font=\small,
      text=cEje, /pgf/number format/fixed, /pgf/number format/precision=1},
    enlarge y limits=0.12]
  \addplot[fill=cApple, draw=none] coordinates {
    (39,MacBook Air M4) (32.5,M6) (27.5,M5)};
  \addplot[fill=cGris!55, draw=none] coordinates {
    (22.2,Snapdragon) (17.8,Intel) (13.4,AMD)};
\end{axis}
\end{tikzpicture}
\caption{Eficiencia multinúcleo de sistemas seleccionados, medida con la
potencia de todo el equipo y una pantalla externa. No mide la duración de la
batería \parencite{notebookcheck2026}.}
\label{fig:eficiencia}
\end{figure}

\section*{3. Efectos más allá de los benchmarks}

La memoria unificada permite que la GPU use la misma memoria que la CPU, sin
el límite de la VRAM de una tarjeta gráfica. Por eso hoy considero al Mac la
mejor opción para desarrollar con modelos de lenguaje locales. La integración
tiene un precio: la memoria se elige al comprar y no se puede ampliar, y
Rosetta tendrá funciones limitadas a partir de macOS~28.

\section*{4. Lo que me llevo como CTO}

Apple afirma que el M6 ofrece hasta 2,4 veces el rendimiento multinúcleo
del M1 \parencite{apple2026m6}, pero la lección no está en el número. A veces basta
con cambiar una pieza; otras, el límite está en cómo trabajan juntas las
piezas. Apple aceptó el costo de rediseñar esas relaciones y construyó un
camino para que las aplicaciones la siguieran.

\printbibliography[title={Referencias}]

\end{document}
